\documentclass[11pt,a4paper]{article}
\usepackage[T1]{fontenc}
\usepackage{lmodern,amsmath,amssymb,amsthm,booktabs,array,microtype}
\usepackage[margin=25mm]{geometry}
\usepackage[hidelinks]{hyperref}
\newtheorem{theorem}{Theorem}
\newtheorem{lemma}[theorem]{Lemma}
\newcommand{\R}{\mathbb R}
\newcommand{\one}{\mathbf 1}
\title{An entropy increase for a structurally disassortative DomiRank graph}
\author{Dev D. Goyal}
\date{26 October 2026}
\usepackage[numbers]{natbib}
\usepackage{xurl}
\begin{document}
\maketitle
\begin{abstract}
A connected graph with 14 vertices and 32 edges disproves the interval
monotonicity assertion in Conjecture 4.9 of Zhang and Zhao,
arXiv:2610.00107v4. Its degree--log-degree edge covariance is strictly
negative. All DomiRank scores are positive on the closed parameter interval
$[251/1000,63/250]$, which lies inside the stable spectral domain.
Nevertheless, the entropy increases between the two endpoints.
The increase exceeds $1/10000$ in natural-logarithm units.
The proof uses exact rational formulas and finite logarithm bounds.
It includes a Lean formalization of the graph, resolvent scores,
interval positivity, covariance, and entropy comparison.
\end{abstract}

\section{Problem and definitions}
Let $G$ be a finite simple undirected graph with adjacency matrix $A$.
Write $d=A\one$ for its degree vector. With dominance threshold one,
the DomiRank score vector at parameter $s$ is
\begin{equation}\label{eq:gamma}
 \gamma(s)=(I+sA)^{-1}(sd).
\end{equation}
For a connected graph with at least two vertices, the stable positive
parameter domain is $0<s<-1/\lambda_{\min}(A)$.
Here $\lambda_{\min}(A)$ is the smallest eigenvalue of the real symmetric
matrix $A$. On an interval where every score is positive, define
\[
 p_i(s)=\frac{\gamma_i(s)}{\sum_j\gamma_j(s)},\qquad
 h(s)=-\sum_i p_i(s)\log p_i(s),\qquad H(s)=\frac{h(s)}{\log 2}.
\]
Thus $H$ is Shannon entropy in bits. The natural logarithm is denoted by
$\log$. Truncating negative scores at zero has no effect on such an interval.

Choose an oriented edge $(U,V)$ uniformly from the $2|E(G)|$ oriented
edges. The structural covariance is
\begin{equation}\label{eq:cov}
 C=\operatorname{Cov}(d_U,\log d_V)
 =\frac{\sum_i(Ad)_i\log d_i}{\sum_i d_i}
 -\frac{\sum_i d_i^2}{(\sum_i d_i)^2}\sum_i d_i\log d_i.
\end{equation}
Changing to base-two logarithms divides $C$ by $\log 2$ and preserves its sign.
The term \emph{structurally disassortative} here means $C\leq0$.

Zhang and Zhao relate this covariance to the initial entropy derivative.
Their Conjecture 4.9 extends the proposed decrease to every interval of
positive scores in the stable domain \cite{zhangzhao}.
The conjecture does not restrict the interval to the first component of
the positive-score set. Their earlier example with positive covariance
does not address this added hypothesis. The graph below has negative
covariance and therefore addresses the refined statement.

\section{The graph and the main result}
Use vertices $0,1,\ldots,13$. The edges are exactly the following unordered pairs:
\begin{itemize}
\item the edge $\{1,2\}$ and the three edges in each of
$\{3,4,5\}$ and $\{6,7,8\}$;
\item every pair with one endpoint in $\{0,1,2\}$ and the other in
$\{3,4,5,6,7,8,9\}$;
\item the four pairs $\{9,j\}$ for $10\leq j\leq13$.
\end{itemize}
These disjoint lists contain $1+6+21+4=32$ edges.
Every vertex is at distance at most two from vertex $9$, so the graph is connected.

\begin{theorem}\label{thm:main}
For this graph, put $s_-=251/1000$ and $s_+=63/250$.
Then $C<0$. For every $s\in[s_-,s_+]$, the matrix $I+sA$ is positive
definite and every coordinate of $\gamma(s)$ is positive. Moreover,
\[
 h(s_-)+\frac1{10000}<h(s_+).
\]
In particular, $H(s_-)<H(s_+)$, contrary to the stated interval monotonicity.
\end{theorem}
The next three sections prove the covariance, stability, and entropy claims.
The construction uses a small leaf probability that increases across the
comparison interval. Exact bounds establish the net entropy increase;
this description alone is not a sign argument.

\section{Covariance and spectral stability}
The degree vector is
\[
 d=(7,8,8,5,5,5,5,5,5,7,1,1,1,1)^T.
\]
Substitution in \eqref{eq:cov} gives
\begin{equation}\label{eq:covvalue}
 C=\frac{159\log5-153\log7-120\log2}{512}<0.
\end{equation}
For the strict sign, use $\log5\leq\log7<3\log2$ and $\log2>0$.
The numerator is at most $6\log7-120\log2$, which is strictly negative.

\begin{lemma}\label{lem:stability}
The matrix $15I+4A$ is positive definite. Consequently $I+sA$ is
positive definite for $0\leq s<4/15$.
\end{lemma}
\begin{proof}
Partition the vertices into classes
\[
 \{0\},\quad\{1,2\},\quad\{3,4,5,6,7,8\},\quad\{9\},\quad\{10,11,12,13\}.
\]
The space of vectors constant on each class is invariant under $A$.
In the basis of class indicators, its quadratic form for $15I+4A$ is
\[
 K=\begin{pmatrix}
15&0&24&4&0\\
0&38&48&8&0\\
24&48&138&0&0\\
4&8&0&15&16\\
0&0&0&16&60
\end{pmatrix}.
\]
Its leading principal determinants are
$15,570,22212,116796,1321488$, all positive. Sylvester's criterion
therefore proves positivity on this five-dimensional space.

Its orthogonal complement has the following invariant decomposition.
The difference on $\{1,2\}$ has adjacency eigenvalue $-1$.
The two sum-zero triangle spaces have eigenvalue $-1$ and total dimension four.
The difference between the two triangle indicators has eigenvalue $2$.
The sum-zero leaf space has eigenvalue $0$ and dimension three.
These dimensions sum to nine. Thus the remaining eigenvalues of $15I+4A$
are $11$ with multiplicity five, $23$ with multiplicity one, and $15$
with multiplicity three. They are positive.
Finally,
\[
 I+sA=\left(1-\frac{15s}{4}\right)I+\frac{s}{4}(15I+4A).
\]
Both summands have nonnegative quadratic forms, and the first is strictly
positive on nonzero vectors when $s<4/15$.
\end{proof}
Since $s_+<4/15$, this proves stability throughout the required interval.
For any eigenpair $Av=\lambda v$ with $v\ne0$ and $\lambda<0$,
positive definiteness gives $1+s\lambda>0$, hence $s<-1/\lambda$.
This applies in particular to the smallest eigenvalue.
The Lean proof checks a full rational factorization of $15I+4A$ instead
of the invariant-space calculation above. Both prove the same positivity statement.

\section{Exact scores on the whole interval}
Put $x=s-1/4$. Then $1/1000\leq x\leq1/500$ on $[s_-,s_+]$.
Define
\begin{align*}
D(x)&=\frac3{16}-\frac{283} {32}x-\frac{59}{4}x^2+52x^3+92x^4+24x^5,\\
N_0(x)&=\frac34-\frac{247}{8}x-\frac{19}{2}x^2+62x^3+24x^4,\\
N_1(x)&=\frac34-\frac{259}{8}x+\frac{13}{2}x^2+86x^3+24x^4,\\
N_2(x)&=\frac14-\frac{123}{8}x+\frac{69}{2}x^2+86x^3+24x^4,\\
N_3(x)&=\frac34-\frac{327}{8}x-\frac{99}{2}x^2+62x^3+24x^4,\\
N_4(x)&=\frac58x+\frac{77}{2}x^2+86x^3+24x^4.
\end{align*}
Each polynomial is positive on this interval. For $D,N_0,N_1,N_2,N_3$,
discard all positive terms except the constant. Bound each remaining
negative term using $x\leq1/500$ and $x^2\leq1/250000$.
The resulting rational lower bounds are positive. Every displayed term of
$N_4$ is nonnegative and its linear term is positive.

The adjacency action on class-constant vectors is the matrix
\[
 B=\begin{pmatrix}
0&0&6&1&0\\0&1&6&1&0\\1&2&2&0&0\\1&2&0&0&4\\0&0&0&1&0
\end{pmatrix}.
\]
Direct polynomial multiplication gives
\[
 (I+sB)(N_0,N_1,N_2,N_3,N_4)^T
   =D(7,8,5,7,1)^T.
\]
Let $c(i)$ be the class index of vertex $i$. Since $I+sA$ is invertible,
\begin{equation}\label{eq:scores}
 \gamma_i(s)=s\frac{N_{c(i)}(s-1/4)}{D(s-1/4)}.
\end{equation}
This proves strict positivity for every real parameter in the interval.
It also establishes the required correspondence with the inverse in
\eqref{eq:gamma}; the rational formulas are not separate score definitions.

\section{A rational entropy certificate}
The common factors $s$ and $D^{-1}$ cancel when the scores are normalized.
Table~\ref{tab:prob} gives the resulting probability of each vertex in a class.
The multiplicities are $(1,2,6,1,4)$.
\begin{table}[htbp]\centering
\begin{tabular}{@{}crr@{}}\toprule
Class & $p_i(s_-)$ & $p_i(s_+)$\\\midrule
0 & $89889445253/534258269542$ & $168020629/988469056$\\
1 & $89703948253/534258269542$ & $167303879/988469056$\\
2 & $29332448253/534258269542$ & $107123383/1976938112$\\
3 & $88634445253/534258269542$ & $81549377/494234528$\\
4 & $82948253/534258269542$ & $685883/1976938112$\\\bottomrule
\end{tabular}
\caption{Exact endpoint probabilities, one value per vertex in each class.}
\label{tab:prob}\end{table}

For $0\leq t<1$, integration of the geometric series gives
\begin{equation}\label{eq:logbound}
 2\sum_{j=0}^{N-1}\frac{t^{2j+1}}{2j+1}
 \leq\log\frac{1+t}{1-t}
 \leq2\sum_{j=0}^{N-1}\frac{t^{2j+1}}{2j+1}
       +\frac{2t^{2N+1}}{1-t^2}.
\end{equation}
The upper remainder bound follows by replacing each remaining denominator
$2j+1$ by $1$ and summing the resulting geometric tail.
Using $t=1/3$ and $N=12$ gives
\[
 \frac{69314718055}{10^{11}}\leq\log2\leq
 \frac{69314718057}{10^{11}}.
\]
For each probability $p$, choose $r$ so that $1\leq2^r p<2$.
Set $t=(2^r p-1)/(2^r p+1)$ and use \eqref{eq:logbound} with $N=7$.
Then subtract $r\log2$. All operations are rational.
Table~\ref{tab:logs} records the resulting enclosures.
\begin{table}[htbp]\centering\small
\begin{tabular}{@{}crrrr@{}}\toprule
&\multicolumn{2}{c}{$10^8\log p_i(s_-)$}&\multicolumn{2}{c}{$10^8\log p_i(s_+)$}\\
Class & lower & upper & lower & upper\\\midrule
0&$-178229885$&$-178229884$&$-177207058$&$-177207057$\\
1&$-178436459$&$-178436458$&$-177634555$&$-177634554$\\
2&$-290218503$&$-290218501$&$-291532324$&$-291532323$\\
3&$-179635882$&$-179635881$&$-180180147$&$-180180146$\\
4&$-877041770$&$-877041769$&$-796635274$&$-796635273$\\\bottomrule
\end{tabular}
\caption{Rational bounds for the logarithms in Table~\ref{tab:prob}.}
\label{tab:logs}\end{table}
Multiplication by the positive probabilities and the class multiplicities gives
\[
 h(s_-)<\frac{2158577}{10^6},\qquad
 \frac{2158712}{10^6}<h(s_+).
\]
Their difference exceeds $135/10^6>1/10000$. This completes the proof
of Theorem~\ref{thm:main}. Dividing by $\log2>0$ proves the entropy increase
in bits. Dividing by $\log14>0$ also proves an increase for normalized entropy.

\section{Formal verification and reproduction}
The Lean sources are in
\path{lean/MathAndTheory/Mathematics/DomiRank/}.
Their namespace is \texttt{MathAndTheory.DomiRank}.
The main theorem is \path{fourteen_vertex_counterexample}.
The formal source correspondence includes the following declarations:
\begin{center}\small
\begin{tabular}{@{}ll@{}}\toprule
Mathematical statement & Lean declaration\\\midrule
Graph adjacency and degrees & \path{A_eq_adjMatrix}, \path{degree_eq_graph_degree}\\
Connectivity & \path{graph_connected}\\
Covariance sign & \path{degreeLogCovarianceBits_neg}\\
Stability on the interval & \path{resolvent_posDef_on_interval}\\
Source spectral condition & \path{spectral_bound_on_interval}\\
Inverse formula and positivity & \path{gamma_eq_candidate}, \path{gamma_pos_on_interval}\\
Normalization correspondence & \path{distribution_eq_theoretical}\\
Strict rational entropy gap & \path{entropyNats_gap}\\
Source-entropy comparison & \path{theoretical_entropy_strict_increase}\\\bottomrule
\end{tabular}\end{center}
The proof contains no numerical eigenvalue assumption. The logarithm
enclosures are finite rational inequalities checked by Lean.
The project pins Lean 4.35.0 and Mathlib commit
\path{56a005855ea17e0b0958f73ec801ae140b0a1679}.
The repository verification record gives the build, axiom audit, and
independent-kernel commands and their actual results.

The independent Python program reconstructs the edges from the three rules
above. It uses exact matrix inversion, rational root isolation for the
whole interval, and an independent rational factorization of $15I+4A$.
It also recomputes the oriented-edge covariance and the finite logarithm
certificate. An 80-decimal numerical check gives
\[
 h(s_+)-h(s_-)\approx0.0001363399223471066814.
\]
This decimal check illustrates the size of the gap. The rational certificate
proves its sign. See \texttt{verification/domirank-entropy/} for the program,
versions, commands, and output.

The theorem concerns the specified mathematical centrality model. It makes
no claim about empirical performance on networks. It refutes the cited
interval statement and makes no claim about a narrower initial-component
version. The literature comparison found no earlier counterexample to
this refined statement; this is not a proof of priority.

\paragraph{Authorship and review.}
Dev D. Goyal directed the AI research and is responsible for publication.
ChatGPT (OpenAI), using the GPT-6 model with Pro reasoning, produced the
supplied construction and proof development. Codex audited the sources,
adapted the formalization, performed independent computations, and prepared
this manuscript. The repository records the later AI referee review.
Machine verification and AI review do not constitute human expert peer review.
This manuscript has not been peer reviewed.
\bibliographystyle{plainnat}
\bibliography{references}
\end{document}
